%=============================================================================
\section{Adversarial Robustness}
\label{sec:adversarial}

\begin{assumption}[Contamination model]
\label{asm:contamination}
A fraction $(1-\delta)$ of judges are honest with competence margin $\gamma > 0$.
A fraction $\delta$ are adversarial.
\end{assumption}

\begin{theorem}[Adversarial breakdown threshold]
\label{thm:adversarial}
For binary votes:
\begin{enumerate}[leftmargin=*,itemsep=2pt]
\item \textbf{Always-wrong adversaries.} If honest judges are correct with
  probability $\frac{1}{2} + \gamma$ and adversaries are always wrong, the
  ensemble is asymptotically correct iff
  \begin{equation}
  \label{eq:adversarial-threshold}
  (1-\delta)\!\left(\tfrac{1}{2} + \gamma\right) > \tfrac{1}{2}
  \quad\Longleftrightarrow\quad
  \delta < \frac{2\gamma}{1 + 2\gamma}.
  \end{equation}

\item \textbf{Sign-flip adversaries.} If adversaries are correct with
  probability $\frac{1}{2} - \gamma$ (i.e., anti-correlated with the truth),
  the threshold becomes
  \begin{equation}
  \label{eq:signflip-threshold}
  \delta < \tfrac{1}{2}.
  \end{equation}
\end{enumerate}
\end{theorem}

\begin{proof}
\textbf{Part 1.} The expected correct-vote fraction is
$(1-\delta)(\frac{1}{2} + \gamma) + \delta \cdot 0
= (1-\delta)(\frac{1}{2} + \gamma)$.
This exceeds $\frac{1}{2}$ iff
$\gamma > \delta(\frac{1}{2} + \gamma)$, i.e.,
$\delta < \gamma / (\frac{1}{2} + \gamma) = 2\gamma/(1+2\gamma)$.

\textbf{Part 2.} The expected correct-vote fraction is
$(1-\delta)(\frac{1}{2}+\gamma) + \delta(\frac{1}{2}-\gamma)
= \frac{1}{2} + \gamma(1-2\delta)$.
This exceeds $\frac{1}{2}$ iff $\gamma(1-2\delta) > 0$, i.e., $\delta < \frac{1}{2}$,
which is the classical Byzantine majority threshold.
\end{proof}

\begin{proposition}[Robust aggregation error]
\label{prop:robust-aggregation}
For bounded scalar scores in $[0,1]$, any robust estimator (e.g., trimmed mean,
median-of-means) inherits an unavoidable $O(\delta)$ bias, giving
\begin{equation}
\mathrm{Error} = O_p(K^{-1/2} + \delta).
\end{equation}
\end{proposition}

\begin{remark}[Optimistic vs.\ pessimistic regimes]
The results in this paper yield a mathematically crisp dichotomy:
\begin{itemize}[leftmargin=*,itemsep=2pt]
\item \textbf{Optimistic}: $\rho_K \to 0$ and $\delta$ below the breakdown
  threshold $\Longrightarrow$ asymptotic consistency.
\item \textbf{Pessimistic}: correlation floor $\rho_0 > 0$ or contamination
  above the honest margin $\Longrightarrow$ irreducible error.
\end{itemize}
\end{remark}
